\textbf{Hardware and software.} One GPU and one CPU model were measured, on one node type. The Tesla K40m is a Kepler-generation GPU (compute capability 3.5, which \texttt{nvcc} 11.7 flags as deprecated), so it is not representative of recent GPUs. Newer GPUs with faster links and more memory bandwidth may move the crossover points; we did not test them. Only g++ 11.4.1, Apple clang 21.0.0 and CUDA 11.7 were used.

\textbf{Workload.} The network is small: its \val{X:weights:GB}\,GB of weight traffic per pass is served from caches, and a larger or convolutional network whose weights do not fit in cache may be limited by memory bandwidth, which would change the CPU--GPU balance. Only float32 inference was measured, with no energy or power measurements.

\textbf{Implementation quality.} The loops and kernels are straightforward: the CPU code is blocked only to allow vectorisation and the CUDA kernel uses coalesced loads without shared-memory tiling; no BLAS library is used. The effective rates at batch 4096 (\val{Sp:serial:gflops:4096}\,GFLOP/s serial, \val{Sp:omp:best:gflops:4096}\,GFLOP/s OpenMP, \val{Sp:cuda:kgflops:4096}\,GFLOP/s for the CUDA kernels alone) show that none of them is tuned. Tuned kernels would change the absolute speedups and possibly the crossover batch sizes; our conclusions concern straightforward implementations. We did not use hardware peak numbers because we did not measure them.

\textbf{Design choices that affect the comparison.} The preprocessing stage is part of our workload definition, and the hybrid gain depends on it; the CUDA implementation deliberately uses one preprocessing thread, which is why the hybrid-sequential control is reported. OpenMP and the pipelined hybrid use the best of seven thread counts per batch size, whereas CUDA (and the hybrid's kernels) use a fixed 16$\times$16 block; the comparison is therefore ``CUDA with a fixed block against OpenMP with tuned thread count'', not best-versus-best. The best block would improve the CUDA end-to-end time by at most \val{Sp:cuda:blockgain:max}\,\% (Section~\ref{sec:block}). One-time GPU initialisation, allocation and weight upload are excluded from the steady-state times, which favours the GPU for short-lived programs; they cost \val{Sp:cold:ctor:mean}\,ms in total (Section~\ref{sec:cold}).

\textbf{Measurement noise.} The Spartan node is managed by a shared scheduler (our jobs held 24 of its 28 CPUs and one of its two GPUs, so other jobs could use the rest), and the M4 ran with background services (system load average about 3 before the runs). Run-to-run variation is small: the median coefficient of variation is \val{Sp:cv:median}\,\% on the Spartan node and \val{Mac:cv:median}\,\% on the M4 (90th percentile \val{Sp:cv:p90}\,\% on Spartan). The largest variation on Spartan is the \val{Sp:slow:impl} configuration at batch \val{Sp:slow:batch} with \val{Sp:slow:threads} threads (\val{Sp:slow:n} passes): mean \val{Sp:slow:mean}\,ms, median \val{Sp:slow:median}\,ms, standard deviation \val{Sp:slow:std}\,ms, minimum \val{Sp:slow:min}\,ms and maximum \val{Sp:slow:max}\,ms, a coefficient of variation of \val{Sp:slow:cv}\,\%. Exactly \val{Sp:slow:count} pass exceeds 1.5 times the median; it is kept in all statistics and in the raw data. On the M4 the largest variation is the serial run at batch \val{Mac:slow:batch}: mean \val{Mac:slow:mean}\,ms, median \val{Mac:slow:median}\,ms, standard deviation \val{Mac:slow:std}\,ms, minimum \val{Mac:slow:min}\,ms and maximum \val{Mac:slow:max}\,ms (\val{Mac:slow:cv}\,\%). The 20 passes of a configuration run back to back, so they are not independent samples of the machine state. macOS ignores thread binding, the Xeon node has two NUMA nodes (threads above 14 cross a socket boundary), and the M4's two core types violate the assumptions of Amdahl's law. The confidence intervals therefore describe the repeatability of these runs, not all possible conditions.

\textbf{Pipelined hybrid.} CUDA events of concurrent streams include time spent waiting for other streams, so the stage durations of the pipelined hybrid are not true durations; we report only its end-to-end time.
